\documentclass[a4paper,11pt]{article}

\usepackage[T1]{fontenc}
\usepackage[utf8]{inputenc}
\usepackage[ngerman]{babel}
\usepackage{lmodern}
\usepackage{microtype}
\usepackage[margin=2.5cm]{geometry}
\usepackage{booktabs}
\usepackage{csquotes}
\usepackage{siunitx}
\sisetup{locale=DE, group-minimum-digits=4}
\usepackage{enumitem}
\usepackage[hyphens]{url}
% Umbruch in langen URLs der Quellenliste an jeder Stelle erlauben
\expandafter\def\expandafter\UrlBreaks\expandafter{\UrlBreaks
  \do\a\do\b\do\c\do\d\do\e\do\f\do\g\do\h\do\i\do\j\do\k\do\l\do\m
  \do\n\do\o\do\p\do\q\do\r\do\s\do\t\do\u\do\v\do\w\do\x\do\y\do\z
  \do\0\do\1\do\2\do\3\do\4\do\5\do\6\do\7\do\8\do\9}
\usepackage[hidelinks]{hyperref}
\hyphenation{OpenCode}

\setlist[description]{style=nextline, leftmargin=1.2em, itemsep=0.3em, parsep=0pt}
\setlength{\parskip}{0.3em}
\setlength{\parindent}{0pt}
\makeatletter
\renewcommand\paragraph{\@startsection{paragraph}{4}{\z@}%
  {1.6ex \@plus .5ex \@minus .2ex}{-1em}{\normalfont\normalsize\bfseries}}
\makeatother

\begin{document}

{\LARGE\bfseries Horizonte für souveräne KI\par}
\smallskip
{\small Stand 17.\,11.\,2026\par}
\medskip
\hrule
\bigskip

%=====================================================================
\section{KI-Horizonte}
\label{sec:horizonte}

Unter \enquote{eigener KI-Infrastruktur} werden an Universitäten sehr
verschiedene Vorhaben und Wünsche zusammengefasst, vom Betrieb eines
offenen Modells für die Verwaltung bis zum Training eigener
Forschungsmodelle.  Sie unterscheiden sich um Größenordnungen in
Kosten, Personal und Realisierbarkeit. Dieser Bericht ordnet sie in
sechs Horizonte und schätzt für jeden ein, wie weit eine deutsche
Universität davon entfernt ist, ihn souverän, also ohne Abhängigkeit
von US-Anbietern und mit eigenen Daten innerhalb der eigenen Organisation, zu erreichen.

\paragraph{Horizont 1: Agentische Computernutzung im Alltag.}
Sprachmodelle erledigen als Agenten alltägliche Wissensarbeit:
Dokumente zusammenfassen, Skripte schreiben, \LaTeX{} korrigieren und
umformatieren, übersetzen, Dokumentendatenbanken durchsuchen (RAG),
als Tutor Schul- und Grundlagenstoff wiederholen. \textbf{Status:
  heute erreichbar.}
 Dieser Horizont sollte unverzüglich umgesetzt
werden.  Die Standard-KI einer Universität muss ein datensparsames, in
der EU oder im eigenen Datacenter gehostetes Modell sein
(Abschnitt~\ref{sec:datenschutz}).  Modelle wie GLM-5.3-Flash
erreichten im Oktober 2026 bei Coding- und Agentenaufgaben das Niveau
des Frontiermodells Claude Opus~4.8 vom Mai 2026
(Abschnitt~\ref{sec:vergleich}).

\paragraph{Horizont 2: Autoformalisierung, Prozessautomatisierung.}
Beispielsweise Übersetzung informeller Mathematik in Lean oder längere
Entwicklungs- und Programmieraufgaben sind \textbf{im Prinzip
  erreichbar}, aber womöglich holpriger als erwartet.  Die
Anforderungen sind mit Horizont~1 vergleichbar, aber diese Modelle
müssen mitunter lange laufen, um (automatisierte) Korrekturschleifen in
Prozesse einzubinden.

\paragraph{Horizont 3: Mathematische Forschung, Beweise finden.}
Frontier-KI findet Beweise für offene Fragen oder trägt zu neuen
Resultaten bei.  \textbf{Status: souverän nicht erreichbar.}  Kein
offenes Modell erreicht Claude Fable~5.2 oder GPT-6.2 Astra.  Das
beste, Kimi~K3.1 Pro, ist an Universitäten in Deutschland nicht für
viele Nutzer hostbar,
lässt sich aber in der EU einkaufen (Abschnitt~\ref{sec:groessen}).
Perspektivisch könnten Verbünde mit kommerziellen Anbietern
kooperieren oder Datacenter in Public-Private-Partnerships betreiben,
in denen auch dieses Hosting möglich wird.

\paragraph{Horizont 4: Eigenes Training.}
Eigene Modelle trainieren, in drei Stufen: Feintuning und Destillation
offener Gewichte, eigenes Pretraining, Frontier-Lauf.  \textbf{Status:
  möglich für sehr kleine Modelle.} Selbst das Training sehr kleiner
Modelle ist ein extrem aufwändiges Unternehmen.  Beispielsweise beschäftigt
NVIDIA etwa 500 Personen und nutzt eine unbekannte Menge an
Hardware/Strom, um ein für Horizont~1 geeignetes Modell selbst zu
trainieren.  Training von Frontiermodellen ist für
Forschungseinrichtungen unrealistisch (Abschnitt~\ref{sec:kosten}).

\paragraph{Horizont 5: Neue Modellarchitekturen.}
Eigene Architekturen, die für die eigene Forschung besser geeignet
sind. \textbf{Status: unrealistisch.} Eine bessere Architektur zu
finden erfordert sehr viele Trainingsläufe, also Rechenzeit und
Personal in der Größenordnung von Horizont~4. Für Mathematik
diskutiert die SAIR-Community dennoch Architektur- und Datenmix-Suche.

\paragraph{Horizont 6: Eigene Spezial-Hardware / Neue KI-Ansätze.}
Eigene Hardware für neue Architekturen, etwa neurosymbolische KI oder
sonstige neuartige Ansätze, die nicht LLM-basiert sind.
\textbf{Status: unklar.}  Es stellen sich sehr viele Fragen.  Selbst
wenn ein Forschungsdurchbruch an einer deutschen Universität passieren
würde, sind die Entwicklung, Vermarktung, Finanzierung oft hier nicht
möglich usw.~usf.  Hier
gibt es beliebig viel Spielraum für Wunschdenken.

\clearpage
%=====================================================================
\section{Querschnittsthemen}
\label{sec:querschnitt}

\subsection{Hardware}

GPU-Hardware ist teuer, 2026 kaum lieferbar und wird schnell
abgeschrieben. Sie muss daher rund um die Uhr ausgelastet sein. Ein
8-GPU-Knoten kostet mehrere hunderttausend Euro. Der Betrieb lässt
sich bei europäischen Anbietern pro Token einkaufen.  Ein eigener
Knoten ist erst bei hoher Auslastung günstiger
(Abschnitt~\ref{sec:hosten}).  \textbf{Selbsthosten wird in der Regel
  nicht billiger, sondern erkauft Datenschutz zu einem Aufpreis.} Das
bestätigt auch die Diskussion in der
SAIR-Community, in der ein Mathematik-Department mit denselben Motiven
(Datenschutz, Unabhängigkeit, fehlende Mittel, vorhandene Hardware)
nach lokalen Modellen fragte.

\subsection{Personal}
Betrieb großer Modelle und Entwicklung von Agentensoftware
(Claude Code, Codex, OpenCode) erfordern Fachkräfte, die in der
Industrie ein Vielfaches einer E11-Stelle verdienen. Auch
professionelle Rechenzentren wie die GWDG betreiben bislang nicht die
größten offenen Modelle.  Universitäten sollten sich an bestehende
Open-Source-Projekte anschließen.  So bauen sie Expertise auf und
halten Anschluss an eine Entwicklung, deren Tempo andere vorgeben.

\subsection{Datenschutz und Politik}
\label{sec:datenschutz}

Bei der KI-Nutzung in Universitäten gibt es zwei Realitäten.
Juristisch gesehen dürfen personenbezogene Daten, die in
Universitäten anfallen (etwa Klausurabgaben, Bewerbungen,
Prüfungsakten), auch bei EU-Cloud-Anbietern nur unter Auflagen
verarbeitet werden: als Auftragsverarbeitung mit Vertrag nach Art.~28
DSGVO.
In der Praxis werden solche Daten nach unserer Beobachtung dennoch
oft über Konten bei US-Anbietern verarbeitet.  Diese müssen Daten
nach dem CLOUD Act auf Anordnung an US-Behörden herausgeben, auch
wenn sie außerhalb der USA gespeichert sind.  Aus internen
Dokumenten, Richtlinien und Senatsberichten entsteht so in
Chatprotokollen ein recht vollständiges Bild der universitären
Abläufe.  Unabhängig von der Absicht der Anbieter ist das eine
politische Abhängigkeit: Was mit diesen Daten geschieht, können die
Universitäten weder kontrollieren noch ausschließen.
Daraus folgt die Priorität für Horizont~1.

%=====================================================================
\section{Fazit}
\label{sec:fazit}
\enlargethispage{\baselineskip}

Horizont~1 und~2 sollten jetzt in Angriff genommen werden, und das
geschieht bereits auf mehreren Ebenen. Rechenzentren wie die GWDG
betreiben offene Modelle für die Wissenschaft. Auch wenn dort noch
nicht die größten Modelle laufen, ist das der richtige
Weg, weil so Betriebserfahrung und Unabhängigkeit entstehen. Auf
europäischer Ebene stellen die AI Factories, in Deutschland HammerHAI
und JAIF, Rechenzeit für die Wissenschaft bereit. Auch die AI
Gigafactories sind richtig, sofern ein Teil der dort garantierten
Rechenzeit der Forschung zufällt (Abschnitt~\ref{sec:projekte}). Vor
Ort kommt es auf Vernetzung an: Universitäten, Rechenzentren und
Fachgruppen sollten Hardware, Modelle, Agenten-Setups und Erfahrungen
teilen, statt jeweils allein anzufangen. Offene Initiativen wie SAIR,
deren erste Phase genau Horizont~1 und~2 umfasst, bieten dafür einen
Rahmen über die eigene Universität hinaus.

Ein Produkt, das alle nutzen, ist dabei noch nicht entstanden: Viele
Nutzer haben bereits ein Abonnement bei einem US-Anbieter, und die
IT-Abteilungen sind ausgelastet.  Mit den Produkten der
Frontier-Anbieter müssen diese Angebote nicht konkurrieren.  Es geht
darum, eigene
Expertise aufzubauen, gerade im Umgang mit Coding-Agenten
(Abschnitt~\ref{sec:agenten}): zu verstehen, was in diesem Feld tatsächlich
geschieht, und Entscheidungen auf Fakten zu stützen statt auf die
Vorstellung, man könne einfach eine eigene Spitzen-KI bauen.

% =====================================================================

\textbf{Im Folgenden gehen wir auf Detailaspekte ein.  Am Ende des Dokuments gibt es ein Glossar, in dem Begriffe wie \enquote{MoE} und andere erklärt werden.} 
\section{Hintergrund: Modellgrößen}
\label{sec:groessen}

\enquote{B} steht für englisch \emph{billion}, also Milliarden
Parameter, \enquote{T} für \emph{trillion}, also Billionen.  Die
Parameter sind die gelernten Zahlen (Gewichte) eines Modells.

Bei einem \emph{dichten} Modell gehen in die Berechnung jedes Tokens
alle Parameter ein.  Fast alle großen Modelle sind dagegen
Mixture-of-Experts-Modelle (MoE).  Ihre Parameter sind auf viele
Teilnetze (Experten) verteilt, und für jedes Token rechnen nur einige
davon.  Deshalb werden für solche Modelle zwei Zahlen angegeben.
\begin{itemize}[itemsep=0.2em]
  \item Die \emph{Gesamtzahl} der Parameter bestimmt den
    Speicherbedarf.  Alle Gewichte müssen im GPU-Speicher liegen, weil
    für jedes Token andere Experten gebraucht werden.
  \item Die \emph{aktiven} Parameter bestimmen den Rechenaufwand pro
    Token und damit Geschwindigkeit, Stromverbrauch und
    Trainingskosten.
\end{itemize}
Die Schreibweise 120B-A12B bedeutet 120~Milliarden Parameter
insgesamt, davon 12~Milliarden aktiv.  GLM-5.3-Flash (320B gesamt,
18B aktiv) braucht also den Speicher eines sehr großen Modells,
rechnet pro Token aber nur so viel wie ein dichtes Modell mit 18B.
Faustregel für den Speicher: 4-Bit-Quantisierung braucht etwa
0,5~GB pro Milliarde Parameter, 8-Bit-Quantisierung 1~GB pro Milliarde
Parameter.  Dazu kommt der Speicher für den Kontext jedes Nutzers.

\subsection{Kleine Modelle für einzelne Nutzer}
\label{sec:klein}

Die Modelle für Smartphone und Laptop sind dicht.  Gesamtzahl und
aktive Parameter fallen dann zusammen.  In der Desktop-Klasse beginnen
die MoE-Modelle.

\begin{itemize}[itemsep=0.2em]
  \item \emph{Smartphone:} ca.\ 2--4B, z.\,B.\ die kleinsten Modelle
    der Reihen Gemma~4 und Qwen3.5.
  \item \emph{Laptop:} ca.\ 27B, z.\,B.
    \href{https://huggingface.co/Qwen/Qwen3.8-27B}{Qwen3.8-27B}. In
    4~Bit quantisiert passt es auf eine Grafikkarte mit 24~GB; auf
    einem Laptop sind ca.\ 32~GB RAM nötig.
  \item \emph{Desktop mit 128~GB RAM:} ca.\ 120B, z.\,B.
    \href{https://build.nvidia.com/nvidia/nemotron-3-super-120b-a12b/modelcard}{Nemotron~3 Super 120B-A12B},
    das weitgehend offen ist, inklusive eines Großteils der
    Trainingsdaten und der Skripte (Abschnitt~\ref{sec:kosten}; Lizenz
    NVIDIA Open Model License bzw.\ OpenMDW). Ein
    Datenpunkt aus der SAIR-Community: Ein DGX Spark (128~GB Unified
    Memory, 273~GB/s, etwa 4\,000--5\,000~USD) mit Qwen~3.8 Flash Next
    (125B gesamt, 6B aktiv, 111~GB in 4~Bit) löst Übungsaufgaben auf
    Hartshorne-Niveau, aber keine Forschungsfragen der algebraischen
    Geometrie.
\end{itemize}

Diese Modelle halluzinieren deutlich und eignen sich vor allem für
abgegrenzte Aufgaben. Mit einem guten Agent-Harness sind sie dennoch
produktiv einsetzbar: \enquote{Lies diese Mail und erstelle den
passenden Kalendereintrag} funktioniert zuverlässig; \enquote{Erstelle
eine Synopse dieser 80 Bewerbungen} gelingt mit einem Coding-Agenten
und passenden \texttt{SKILL.md}-Dateien.

\subsection{Workstation- und Datacenter-Klasse}

Hier sind mehrere GPUs mit eigenem HBM-Speicher nötig, die per NVLink
verbunden sind. Für Einzelnutzer genügt teils ein Rechner mit sehr viel
Unified Memory.

\begin{table}[htbp]
  \centering
  \begin{tabular}{lll}
    \toprule
    Modell & Parameter gesamt / aktiv & Gewichte in 4~Bit ca. \\
    \midrule
    \href{https://huggingface.co/zai-org/GLM-5.3-Flash}{GLM-5.3-Flash} & 320B / 18B & 160~GB \\
    \href{https://huggingface.co/deepseek-ai/DeepSeek-V4.1-Flash}{DeepSeek-V4.1-Flash} & 552B + 196B Memory / 8--16B & 380~GB \\
    \href{https://huggingface.co/zai-org/GLM-5.3}{GLM-5.3} & 744B / 40B & 375~GB \\
    \href{https://huggingface.co/moonshotai/Kimi-K3}{Kimi~K3} & 2,8T / 16 von 896 Experten & 1,4~TB \\
    \bottomrule
  \end{tabular}
  \caption{Open-Weights-Modelle der Datacenter-Klasse (Stand Oktober 2026).}
  \label{tab:modelle}
\end{table}

Diese Modelle sind für sehr viele Einsatzzwecke ausreichend. Kimi~K3
lag beim Launch in Benchmarks nur hinter Fable und GPT-5.6 Sol.  Die
aktuelle Frontier (Claude Fable~5.2, GPT-6.2 Astra, Gemini~4 Carbon,
Stand 17.\,11.\,2026) ist eine Generation weiter.  Seit dem
23.\,10.\,2026 gibt es mit Kimi~K3.1 Pro (2,8T gesamt, 104B aktiv)
einen Nachfolger, ebenfalls mit offenen Gewichten.
Die Gewichte aller genannten Modelle sind frei herunterladbar, die
Lizenzen unterscheiden sich (MIT, Apache~2.0, bei Kimi~K3 eine eigene
Lizenz). Der Betrieb ist teuer: Für Kimi~K3 sind allein für die
Gewichte 1,4~TB GPU-Speicher nötig, dazu der KV-Cache für den Kontext
jedes Nutzers, der bei langen Kontexten (1M Token) und vielen
gleichzeitigen Nutzern mehr Speicher brauchen kann als die Gewichte.
Die Größe der Frontiermodelle ist nicht öffentlich;
unbestätigten Berichten zufolge liegt Fable bei etwa 10T Parametern,
Astra darüber.

\subsection{Leistungsvergleich}
\label{sec:vergleich}

Es gibt sehr viele Benchmarks für Sprachmodelle, für Wissen,
Mathematik, Programmieren und agentische Aufgaben, und jede Woche
kommen neue Modelle hinzu. Einen guten Überblick bietet Artificial
Analysis (\url{https://artificialanalysis.ai}): Die Seite misst
offene und geschlossene Modelle mit denselben Benchmarks und stellt
Leistung, Preis und Geschwindigkeit nebeneinander.

Tabelle~\ref{tab:vergleich} vergleicht offene und geschlossene Modelle
nach den Messungen von Artificial Analysis, die unabhängig von den
Herstellern erhoben werden. Es sind Benchmark-Ergebnisse, also
berichtet, nicht im eigenen Einsatz verifiziert. Der Intelligence
Index fasst zehn Evaluationen zusammen, davon vier agentische;
Terminal-Bench~4.0 und GDPval-AA sind rein agentisch. GLM-5.3-Flash
liegt in dieser Messung gleichauf mit Claude Opus~4.8 und bei den
agentischen Benchmarks meist darüber. Das damals stärkste offene
Modell, GLM-5.3, lag rund 13 Indexpunkte hinter dem damals besten
Modell, Claude Opus~5.5.

\begin{table}[htbp]
  \centering
  \footnotesize
  \begin{tabular}{llrrr}
    \toprule
    Modell & Release & Intelligence Index & Terminal-Bench 4.0 & GDPval-AA (Elo) \\
    \midrule
    GLM-5.3-Flash (offen) & 08/2026 & 41,8 & 32,8~\% & \num{1644} \\
    GLM-5.3 (offen) & 08/2026 & 44,8 & 41,9~\% & \num{1650} \\
    Kimi~K3 (offen) & 07/2026 & 43,6 & 12,6~\% & \num{1533} \\
    Claude Opus~4.6 & 02/2026 & 31,9 & -- & -- \\
    Claude Opus~4.8 & 05/2026 & 41,8 & 21,7~\% & \num{1453} \\
    Claude Fable~5.1 & 09/2026 & 53,4 & 52,0~\% & \num{1754} \\
    GPT-6 Astra & 09/2026 & 52,7 & 59,1~\% & \num{1573} \\
    Claude Opus~5.5 & 09/2026 & 57,6 & 59,6~\% & \num{1866} \\
    \bottomrule
  \end{tabular}
  \caption{Leistungsvergleich nach Artificial Analysis (Intelligence
    Index v4.3.2, Stand 16.\,10.\,2026). Jeweils höchste
    Reasoning-Stufe; Opus~4.6 ohne Werte auf den aktuellen agentischen
    Benchmarks.}
  \label{tab:vergleich}
\end{table}

Tabelle~\ref{tab:vergleich} gibt den Stand vom 16.\,10.\,2026 wieder.
Seitdem hat Artificial Analysis den Index neu zusammengesetzt
(Version~5.0.1, neun Evaluationen, davon sechs mit Werkzeugnutzung).
Die neuen Werte sind mit den älteren nicht vergleichbar, und Claude
Opus~4.8 wird nicht mehr bewertet.  Stand 17.\,11.\,2026 führt Claude
Fable~5.2 mit 65~Punkten vor GPT-6.2 Astra (62), Gemini~4 Carbon (60)
und Claude Opus~5.5 (59).  Das stärkste offene Modell ist Kimi~K3.1
Pro mit 46~Punkten, GLM-5.3 erreicht 32 und GLM-5.3-Flash 28.
Ranglisten dieser Art ändern sich innerhalb von Wochen.  Jede
Einordnung eines Modells ist deshalb mit ihrem Datum zu lesen.

\subsection{Agenten und Korrekturschleifen}
\label{sec:agenten}

Benchmarks messen vor allem das Modell. Im Einsatz entscheidet aber
das Gesamtsystem aus Modell, Harness und Werkzeugen. Viele Schwächen
der Modelle, allen voran Halluzinationen, werden in der Praxis nicht
durch bessere Modelle behoben, sondern durch Korrekturschleifen: Ein
Coding-Agent führt Code aus, lässt Tests laufen, liest Fehlermeldungen
und korrigiert sich. In der Mathematik übernimmt der Lean-Kernel diese
Rolle. Ein offenes Modell in einem guten Agenten leistet deshalb
deutlich mehr, als seine Benchmark-Werte vermuten lassen. \textbf{Die
  Auseinandersetzung mit Coding-Agenten, ihren Harnesses und
  Arbeitsweisen ist Grundlagenarbeit und mindestens so wichtig wie die
  Wahl des Modells.}

\subsection{Selbst hosten oder einkaufen: eine Überschlagsrechnung}
\label{sec:hosten}

Lohnt sich ein eigener GPU-Knoten gegenüber dem Einkauf pro Token?
Wir rechnen das für GLM-5.3-Flash durch, das Modell aus Horizont~1.
Alle Eingangsgrößen sind berichtet, nicht selbst gemessen, und stammen
vom November 2026.

\emph{Eigener Knoten.}  Die Gewichte in 8~Bit belegen 306~GiB und
passen damit nicht auf zwei H200 (zusammen 282~GB).  Das Standardrezept
der Serving-Software vLLM sieht acht H200 vor.  Die 4-Bit-Fassung
setzt laut Rezept die neuere GPU-Generation Blackwell voraus.  Ein
Server mit acht H200 kostet netto rund \num{318000}~USD und nimmt
9,1~kW auf.  Über vier Jahre abgeschrieben sind das rund
\num{80000}~USD pro Jahr.  Dazu kommen rund \num{16000}~EUR Strom
(PUE~1,2 und 17,2~ct/kWh), zusammen knapp \num{100000} pro Jahr.
Zwischen USD und EUR unterscheiden wir bei dieser Genauigkeit nicht.
Nicht enthalten sind Personal, Raum, Netz und Ausfälle.

\emph{Durchsatz.}  Eine veröffentlichte Messung auf dieser Hardware
(acht H200, vLLM, Anfragen mit \num{8192}~Tokens Eingabe und
256~Tokens Ausgabe) ergibt bei 128 gleichzeitigen Anfragen
\num{2359}~Ausgabe-Tokens und 9,3~Anfragen pro Sekunde.  Rund um die
Uhr voll ausgelastet verarbeitet der Knoten damit im Jahr 2,4~Bio.\
Eingabe- und 74~Mrd.\ Ausgabe-Tokens.

\emph{Einkauf.}  Cortecs, ein Vermittler für Anbieter mit Sitz in der
EU, verlangt für GLM-5.3-Flash 0,09~EUR pro Million Eingabe-Tokens und
0,314~EUR pro Million Ausgabe-Tokens.  Die genannte Tokenmenge kostet
dort rund \num{240000}~EUR.  Beim US-Vermittler OpenRouter reichen die
Preise der 44 Endpunkte für dieses Modell von 0,04 und 0,08~USD bis
0,23 und 0,75~USD pro Million Eingabe- und Ausgabe-Tokens.  Der
Hersteller Z.ai verlangt 0,15 und 0,50~USD, der einzige dort als
europäisch ausgewiesene Anbieter 0,175 und 0,49~USD.

\emph{Ergebnis.}  Bei Vollauslastung kosten eingekaufte Tokens etwa
das Zweieinhalbfache der Hardware- und Stromkosten des eigenen
Knotens.  Der eigene Knoten ist also ab etwa 40~Prozent
Dauerauslastung günstiger.  Rechnet man eine Stelle für den Betrieb
hinzu (Annahme: \num{90000}~EUR pro Jahr), liegt die Schwelle bei
knapp 80~Prozent.  Eine Arbeitswoche von 40 bis 50~Stunden entspricht
24 bis 30~Prozent der Woche.  Eine einzelne Universität erreicht die
Schwelle daher kaum, ein Verbund oder ein Rechenzentrum für viele
Einrichtungen eher.  Die Rechnung ist grob.  Sie gilt nur für das
genannte Lastprofil, und Prompt-Caching, das beim Einkauf den Preis
und auf dem eigenen Knoten die Last senkt, ist nicht berücksichtigt.

%=====================================================================
\section{Trainingskosten}
\label{sec:kosten}

Dieser Abschnitt unterlegt Horizont~4 mit öffentlich zugänglichen
Zahlen.  Für Frontiermodelle sind Kosten und Rechenaufwand fast durchweg geheim. Kurzfassung: Der
finale Trainingslauf eines Modells bis in die 100B-MoE-Klasse
(\enquote{Desktop}) kostet Millionen. Die großen Posten sind Personal,
Datenarbeit und Hardwarezugang.\footnote{Quellen in
  Abschnitt~\ref{sec:quellen}.}

Ein Training umfasst mehrere Stufen: Vorab werden Daten gesammelt,
gefiltert und in kleineren Ablationsläufen getestet, dann folgt das
Pretraining auf Billionen Tokens als teuerster Rechenschritt. Danach
kommen Mid- und Long-Context-Training sowie das Post-Training aus
überwachtem Feintuning (SFT) und Reinforcement Learning, das etwa
Reasoning und Werkzeugnutzung beibringt.

\subsection{Offengelegte Zahlen}

Nur wenige Entwickler nennen Geldbeträge. DeepSeek bezifferte den
finalen Lauf von DeepSeek-V3 auf 2,788~Mio.\ H800-Stunden, also
5,6~Mio.\ USD bei 2~USD pro GPU-Stunde, ausdrücklich ohne
Vorexperimente und ohne Hardwarekauf. Ai2 nannte für Olmo~3 32B
1,38~Mio.\ H100-Stunden und 2,75~Mio.\ USD. Meta nennt GPU-Stunden,
aber kein Geld. Alibaba legt für Qwen die Tokenzahlen (also
verarbeitete Token beim Training) offen (18T für Qwen2.5, 36T für
Qwen3), aber in keiner Generation GPU-Typ, GPU-Zahl oder
Kosten.  Das größte Modell der nächsten Generation, Qwen4-Max, wird
seit November 2026 über Alibabas API angeboten, bislang ohne offene
Gewichte (berichtet, Stand 17.\,11.\,2026).  NVIDIA nennt für
Nemotron~3 weder GPU-Stunden noch
Kosten. Tabelle~\ref{tab:kosten} fasst zusammen und trennt dabei, was
die Beträge enthalten: Die meisten sind der Mietwert der GPU-Stunden
des finalen Laufs, einige sind Schätzungen von Epoch AI, andere
Gesamtbudgets ganzer Programme. Geschätzte GPU-Stunden sind aus der
Formel $C = 6\,N\,T$ (Rechenaufwand in FLOP gleich sechsmal aktive
Parameter mal Tokens) abgeleitet.

\begin{table}[htbp]
  \centering
  \footnotesize
  \setlength{\tabcolsep}{3pt}
  \begin{tabular}{llllll}
    \toprule
    Modell (Datum) & Parameter & Tokens & GPU-Stunden & Kosten & Basis \\
    \midrule
    DeepSeek-V3 (12/2024) & 671B / 37B & 14,8T & 2,79 Mio.\ H800 & 5,6 Mio.\ USD & Miete, D \\
    Olmo 3 32B (11/2025) & 32B & 5,9T & 1,38 Mio.\ H100 & 2,75 Mio.\ USD & Miete, D \\
    Llama 3.1 405B (07/2024) & 405B & 15T & 30,8 Mio.\ H100 & 62 Mio.\ USD & Miete, D \\
    Llama 4 Maverick (04/2025) & 400B / 17B & 22T & 2,4 Mio.\ H100 & ca.\ 5 Mio.\ USD & Miete, D \\
    gpt-oss-120b (08/2025) & 117B / 5B & k.\,A. & 2,1 Mio.\ H100 & ca.\ 4 Mio.\ USD & Miete, D \\
    Nemotron 3 Super (03/2026) & 121B / 12B & 25T & k.\,A. & 1--2 Mio.\ USD & Miete, E \\
    Qwen3-32B (05/2025) & 33B & 36T & ca.\ 5 Mio.\ H100 & 10--15 Mio.\ USD & Miete, E \\
    Qwen3-235B-A22B (05/2025) & 235B / 22B & 36T & ca.\ 4,4 Mio.\ H100 & 9--13 Mio.\ USD & Miete, E \\
    \midrule
    Nemotron-4 340B (06/2024) & 340B & 9T & 12--13 Mio.\ H100 & 21 Mio.\ USD & Epoch \\
    Grok 4 (07/2025) & k.\,A. & k.\,A. & 246 Mio.\ H100 & 490 Mio.\ USD & Epoch \\
    \midrule
    BLOOM 176B (07/2022) & 176B & 0,35T & 1,08 Mio.\ A100 & 3--5 Mio.\ EUR & Budget \\
    Apertus-70B (09/2025) & 71B & 15T & 6--9 Mio.\ GH200 & CHF 20 Mio. & Budget \\
    Teuken-7B (11/2024) & 7B & 4--6T & k.\,A. & 14 Mio.\ EUR & Budget \\
    \bottomrule
  \end{tabular}
  \caption{Offengelegte und geschätzte Trainingskosten. Parameter:
    gesamt / aktiv bei MoE, sonst dichte Modelle. \emph{Miete:}
    GPU-Stunden des finalen Laufs mal 2~USD, ohne Vorexperimente,
    Daten, Personal und Anschaffung; D: GPU-Stunden vom Hersteller
    offengelegt, E: geschätzt. \emph{Epoch:} Schätzung von Epoch AI
    (bei Nemotron-4 amortisierte Hardware und Energie); für Llama~3.1
    schätzt Epoch 170~Mio.\ USD (AI Index 2025). \emph{Budget:}
    Gesamtbudget eines Programms einschließlich Personal (Apertus:
    Swiss AI Initiative, Teuken: OpenGPT-X-Förderung, BLOOM: inkl.\
    Engineering laut Sekundärquelle), mit den übrigen Zeilen nicht
    vergleichbar. Quellen: technische Berichte und Modellkarten der
    Hersteller, Epoch AI.}
  \label{tab:kosten}
\end{table}

\textbf{Für eine Budgetplanung taugen diese Zahlen nicht, für
  Größenordnungen und Vergleiche schon.} Die Mietwerte sind
Untergrenzen: Sie enthalten weder die vorgelagerten Experimente, für
die das Rechenbeispiel unten noch einmal das 1,5- bis 2-Fache ansetzt,
noch Datenarbeit und Personal, das bei Frontiermodellen nach Epoch AI
29--49~Prozent der Gesamtkosten ausmacht. Wie groß die Lücke sein
kann, zeigt OpenGPT-X (Abschnitt~\ref{sec:dsp}): 14~Mio.\ EUR Förderung
für ein 7B-Modell, dessen finaler Lauf rechnerisch unter 1~Mio.\ EUR
Mietwert liegt.

Drei Beobachtungen. Erstens liegt zwischen einem offenen Lauf
(2,75--5,6~Mio.\ USD) und einem Frontier-Lauf (100--500~Mio.\ USD)
etwa Faktor~100. Zweitens beginnt niemand bei null: NVIDIAs
\enquote{vollständig offene} Nemotron-Daten wurden mit Qwen- und
DeepSeek-Modellen synthetisch angereichert, und NVIDIAs Dokumentation
beziffert die freigegebenen Pretraining-Daten auf 40 bis 50~Prozent
des internen Mixes. Qwen3 liefert die eindrücklichste Zahl zum Preis
des Neuanfangs: Das Reasoning-Post-Training von Qwen3-8B kostete
\num{17920}~GPU-Stunden per Reinforcement Learning, aber nur
\num{1800}~GPU-Stunden per Destillation aus dem größeren Modell. Das
Pretraining desselben 8B-Modells liegt rechnerisch bei 1,25~Mio.\
GPU-Stunden. Für die RL-Stufe von DeepSeek-R1 nennt DeepSeek rund
\num{294000}~USD.  Drittens lässt sich ein offener Lauf dieser Klasse
auch außerhalb der großen Firmen organisieren.  Das Projekt Marin
(Abschnitt~\ref{sec:projekte}) hat am 5.\,11.\,2026 das Pretraining
eines MoE-Modells mit 535B Parametern, davon 23B aktiv, auf 18T
Tokens abgeschlossen.  Die 768 GPUs vom Typ GB200 stellte eine
Stiftung bereit.  Den Rechenaufwand beziffert das Projekt auf rund
$2{,}6\cdot 10^{24}$~FLOP.  Mit den Annahmen des folgenden
Rechenbeispiels entspricht das etwa 2,4~Mio.\ H100-Stunden oder
7~Mio.\ EUR Mietwert.

\subsection{Rechenbeispiel für universitätsrelevante Größen}

Tabelle~\ref{tab:rechenbeispiel} rechnet Trainingsläufe in den Größen
durch, die für ein deutsches Konsortium in Frage kommen. Annahmen:
$C = 6NT$ mit $N$ als aktive Parameter; H100 mit 989~TFLOP/s
Spitzenleistung; Auslastung (MFU) 40~Prozent für dichte Modelle
(Llama~3 erreichte 38--43~Prozent) und 30~Prozent für MoE; Miete
3~EUR pro H100-Stunde (Scaleway 2,87~EUR, OVHcloud 3,33~EUR, Stand
Herbst 2026); Energie 1,28~kWh pro H100-Stunde inklusive Knoten, Netz
und Kühlung bei PUE 1,2; Strompreis 17,2~ct/kWh (BDEW 2026,
Mittelstand). Nicht enthalten sind Ablationen, Neustarts,
Post-Training, Datenaufbereitung und Personal; dafür ist das 1,5- bis
2-Fache des finalen Laufs zu budgetieren.

\begin{table}[htbp]
  \centering
  \footnotesize
  \setlength{\tabcolsep}{4pt}
  \begin{tabular}{lrlrrr}
    \toprule
    Modell / Tokens & H100-h & Cluster (3 Monate) & Miete & Energie & Strom (DE) \\
    \midrule
    7B dicht / 2T & \num{59000} & 27 GPUs & 0,18 Mio.\ EUR & 76 MWh & \num{13000} EUR \\
    7B dicht / 4T & \num{118000} & 55 GPUs & 0,35 Mio.\ EUR & 151 MWh & \num{26000} EUR \\
    30B dicht / 10T & 1,27 Mio. & 590 GPUs & 3,8 Mio.\ EUR & 1,6 GWh & 0,28 Mio.\ EUR \\
    120B-A12B MoE / 15T & 1,0 Mio. & 470 GPUs & 3,0 Mio.\ EUR & 1,3 GWh & 0,22 Mio.\ EUR \\
    120B-A12B MoE / 25T & 1,7 Mio. & 780 GPUs & 5,0 Mio.\ EUR & 2,2 GWh & 0,37 Mio.\ EUR \\
    400B dicht / 15T & 25 Mio. & \num{11700} GPUs & 76 Mio.\ EUR & 32 GWh & 5,6 Mio.\ EUR \\
    \bottomrule
  \end{tabular}
  \caption{Rechenbeispiel für finale Trainingsläufe, eigene Arithmetik
    mit den im Text genannten Annahmen. Die Zeile 120B-A12B/25T
    entspricht dem Rezept von Nemotron~3 Super, die Zeile 400B/15T
    der Klasse von Llama~3.1 405B.}
  \label{tab:rechenbeispiel}
\end{table}

Die Plausibilitätsprüfung gegen gemessene Werte fällt gut aus: OLMo~2
7B maß 157~MWh für 4T Tokens; Meta nannte 30,8~Mio.\ H100-Stunden für
Llama~3.1 405B einschließlich Neustarts und Langkontextphasen.

Auf JUPITER übersetzt: Ein einziges GCS-Large-Scale-Projekt an der
Schwelle von 1,73~Mio.\ GPU-Stunden pro Jahr deckt das 30B/10T- oder
das 120B-MoE/15T-Beispiel ab. Ein dichter 400B-Lauf entspräche mit
25~Mio.\ GPU-Stunden fast dem gesamten Angebot einer
GCS-Ausschreibung von 27,3~Mio.\ GPU-Stunden (Stand Juli 2026), für
die das GCS rund 90~Anträge erwartet. Die
realistische Pretraining-Obergrenze für ein Konsortium liegt damit bei
10--30B dichten oder etwa 120B-A12B MoE-Modellen pro Jahr. Die
Apertus-Klasse (70B auf 15T Tokens, 6--9~Mio.\ GPU-Stunden) war in der
Schweiz nur möglich, weil 40~Prozent eines Exascale-Systems für ein
Quartal reserviert wurden.

\subsection{Daten, Strom, Personal}
\label{sec:dsp}

\emph{Daten.} Die Rohkorpora sind frei: FineWeb (15T Tokens), DCLM
(4T), Nemotron-Pretraining-Dataset (6,6T, davon rund 3T synthetisch),
The Stack v2 (900B Tokens Code). Teuer ist die Iteration: Allein die
Ablationen für FineWeb kosteten \num{80000}~H100-Stunden, und 1T Tokens
synthetische Daten über eine API kosten aktuell 0,75 bis 1,5~Mio.\
USD. Deutscher Text ist zwei Größenordnungen knapper als englischer:
FineWeb-2 enthält 262~Mrd.\ deutsche Wörter, die lizenzrechtlich
sauberen \enquote{German Commons} 155~Mrd.\ Tokens.

\emph{Recht.} Rechtlich stehen Hochschulen gut da.  \S\,60d UrhG
erlaubt Forschungsorganisationen Text- und Data-Mining für die
wissenschaftliche Forschung.  Einen Nutzungsvorbehalt der
Rechteinhaber sieht das Gesetz nur für die allgemeine Schranke in
\S\,44b UrhG vor, nicht in \S\,60d.  Das OLG Hamburg hat im
LAION-Verfahren am 10.\,12.\,2025 entschieden, dass sich
nichtkommerzielle Forschungsorganisationen auf \S\,60d stützen
können, und die Revision zum BGH zugelassen.  Das LG München~I hat
allerdings im Verfahren GEMA gegen OpenAI (11.\,11.\,2025) die
Memorisierung von Liedtexten in Modellgewichten als Vervielfältigung
gewertet.  Dieses Urteil ist nicht rechtskräftig.  Der AI Act
verlangt von Anbietern von Modellen mit allgemeinem Verwendungszweck
unter anderem eine Strategie zur Einhaltung des Urheberrechts und eine
öffentliche Zusammenfassung der Trainingsinhalte (Art.~53).  Als
Anhaltspunkt für ein solches Modell nennen die Leitlinien der
Kommission einen Trainingsaufwand von mehr als $10^{23}$~FLOP.  Diesen
Wert überschreitet bereits das 7B/4T-Beispiel in
Tabelle~\ref{tab:rechenbeispiel}.  Die Pflichten greifen, sobald ein
Modell in der EU in Verkehr gebracht wird, wozu die Leitlinien auch
das Hochladen in ein öffentliches Repositorium zählen.  Forschung und
Entwicklung vor diesem Schritt sind ausgenommen, und für Modelle unter
freier Lizenz entfällt ein Teil der Dokumentationspflichten.

\emph{Aufgaben statt Texte.} Für das Pretraining gibt es den
beschriebenen Fundus frei verfügbarer Texte, für das Post-Training
nichts Vergleichbares.  Open Athena fasst das so: Pretraining
verbraucht Dokumente, Post-Training verbraucht Aufgaben, also eine
Umgebung, ein Ziel, Werkzeuge und eine automatische
Bewertung.\footnote{D.~Hall, \emph{There Is No Post-Training Data
  Commons}, 15.\,10.\,2026,
  \url{https://openathena.ai/blog/there-is-no-posttraining-data-commons/}.
  Die Zahlen zu SWE-bench sind von dort übernommen.}  Selbst die
Softwareentwicklung, die ihre Arbeit in Versionsverwaltung und Tests
weitgehend aufzeichnet, liefert wenig: Für den Benchmark SWE-bench
blieben von gut \num{90000} Pull Requests \num{2294}~Aufgaben übrig.
OpenAI ließ danach 93~erfahrene Python-Entwickler \num{1699} dieser
Aufgaben von Hand prüfen.  Mehr als zwei Drittel wurden beanstandet,
500 bildeten am Ende SWE-bench Verified.  Die großen Firmen bezahlen
deshalb Fachleute dafür, ihre Arbeit in solche Aufgaben zu
übersetzen.  Für Universitäten folgt daraus zweierlei.  Wie Fachleute
in einer konkreten Situation vorgehen, ist ein wertvolles Gut.  Es
entsteht beim Arbeiten mit Agenten nebenbei und fällt bei dem
Anbieter an, dessen Dienst genutzt wird.  Und die Mathematik ist
womöglich in einer günstigeren Lage als andere Fächer, denn eine
formalisierte Aussage mit dem Lean-Kernel als Prüfer ist bereits eine
Aufgabe mit automatischer Bewertung.  Ob die vorhandenen formalen
Bibliotheken dafür ausreichen, ist offen.

\emph{Strom.} In Deutschland kostet die Energie für eine H100-Stunde
0,19 bis 0,22~EUR, in Finnland ein Drittel davon. Das sind rund
10~Prozent der Mietkosten, aber für eine Universität mit
\emph{gewährter} Rechenzeit der einzige Posten, der aus dem Haushalt
bezahlt wird: \num{13000}~EUR für ein 7B-Modell, \num{280000}~EUR für
ein 30B-Modell. Das Energieeffizienzgesetz erfasst Rechenzentren ab
300~kW Anschlussleistung.  Wer ab Juli 2026 den Betrieb aufnimmt, muss
eine PUE von höchstens 1,2 erreichen und mindestens 10~Prozent der
Energie als Abwärme wiederverwenden.

\emph{Personal.} Was die Tabellen nicht zeigen, ist meist größer als
das, was sie zeigen. OpenGPT-X erhielt rund 14~Mio.\ EUR für zehn
Partner über drei Jahre und lieferte ein 7B-Modell, dessen reiner
Trainingslauf rechnerisch unter 1~Mio.\ EUR Mietwert liegt. Apertus
hatte ein Team von über 100 Personen, Nemotron rund 500
Vollzeitäquivalente, Marin (Abschnitt~\ref{sec:projekte}) zehn
Vollzeit-Engineers. Epoch AI setzt Personal auf 29--49~Prozent der
Gesamtkosten eines Frontiermodells an, Energie auf 2--6~Prozent.
Zur Qualität: Teuken-7B erreicht über 21 EU-Sprachen 0,543 gegenüber
0,563 für Llama-3.1-8B, Apertus-70B 67,5~Prozent gegenüber 69,8~Prozent
für Qwen2.5-72B. Europäische öffentliche Modelle erreichen damit fast
das Niveau älterer offener Modelle.

\subsection{Drei Stufen}

Aus den Zahlen folgt für Horizont~4 eine Staffelung:

\begin{enumerate}[itemsep=0.2em]
  \item \emph{Feintuning und Destillation offener Gewichte} (Qwen,
    Nemotron-Basis-Checkpoints, Apertus): Tausende GPU-Stunden, für
    jede Universität mit einem GPU-Knoten erreichbar. Das ist die
    Rechenseite von Horizont~1 und~2 und der realistische Kern der
    SAIR-Initiative.
  \item \emph{Eigenes Pretraining im Bereich 7--30B} auf offenen Korpora
    unter \S\,60d UrhG: 0,1 bis 1,3~Mio.\ H100-Stunden und
    \num{13000} bis \num{280000}~EUR Strom, über GCS oder EuroHPC
    finanzierbar. Das Ergebnis ist ein Modell der Laptop-Klasse
    (Abschnitt~\ref{sec:klein}) auf dem Stand von vor ein bis zwei
    Jahren.
  \item \emph{Frontier-Lauf}: 25~Mio.\ GPU-Stunden, 32~GWh, 5--6~Mio.\
    EUR Strom, dazu ein Team von 50--100 Personen über Jahre. Der
    Engpass ist nicht die Rechenzeit, sondern dieses Team. In
    Deutschland ist das nur als nationales oder europäisches Vorhaben
    denkbar.  US-Firmen setzen 100--500~Mio.\ USD pro Lauf ein.
\end{enumerate}

\textbf{Der Ertrag von Stufe~2 ist die Expertise, nicht das Modell.}
Offene Modelle der Laptop-Klasse, etwa Qwen3.8-27B, stehen bei Hugging
Face zum Herunterladen bereit und sind einem ersten eigenen Lauf
voraus. Das Stromkosten-Delta von Faktor~2--3
gegenüber Skandinavien ist eine Standortfrage, die zu benennen ist.
Stufe~3 sollte nicht versprochen werden. Auf Stufe~1 aufbauend lässt
sich prüfen, ob eigenes Post-Training offene Modelle in Mathematik
verbessert.  Gemeint sind Feintuning auf Beispielen und Reinforcement
Learning, bei dem das Modell an automatisch bewerteten Aufgaben lernt
(Abschnitt~\ref{sec:dsp}).  Der Engpass sind dabei weniger die
Rechenzeit als die Aufgaben und die Menschen, die sie aufbereiten.

%=====================================================================
\section{Projekte und Initiativen}
\label{sec:projekte}

Stand 16.\,10.\,2026, nach Rolle sortiert, ohne Anspruch auf
Vollständigkeit.  Die Einträge zu Marin, Apertus, OpenEuroLLM und den
AI Gigafactories sind am 17.\,11.\,2026 aktualisiert.

\subsection{Offene Modelle für Mathematik}

\begin{description}
\item[SAIR Open Math Model Initiative.]
  Foundation for Science and AI Research, 2025 von Terence Tao mit
  privaten Spendern gegründet; Förderer der Wettbewerbe sind XTX
  Markets und Susquehanna. Am 18.\,09.\,2026
  angekündigt:\footnote{\url{https://terrytao.wordpress.com/2026/09/18/sairs-open-math-model-initiative/}
  und \url{https://sair.foundation/open-math-model/}, abgerufen am
  10.\,10.\,2026.} offene Modelle und offenes Tooling für die
  Mathematik. Die erste Phase sind Alltagswerkzeuge (Argumente
  verstehen, Referenzen prüfen, Beispiele explorieren, Code schreiben,
  Beweise formalisieren); sie entspricht genau Horizont~1 und~2. Es
  geht um Unabhängigkeit der Community von den großen KI-Firmen,
  Reproduzierbarkeit und Governance, nicht um geografische
  Souveränität. Partner, Plan, Modellgrößen und Budget sind noch nicht
  benannt (Stand 10.\,10.\,2026). Offen ist die Lizenz: Die Ankündigung
  nennt Apache/MIT/CC~BY, in der Community werden AGPLv3 und CC~BY-SA
  vorgeschlagen, damit KI-Firmen Community-Arbeit nicht ohne
  Gegenleistung einbauen. Unterstützung angeboten haben unter anderem
  Hugging Face (Trainings-Know-how, Daten, Compute; Betreiber von
  FineWeb und SmolLM), die Open Model Initiative der Linux Foundation
  und Open Athena.
\item[Open Athena und Marin (Jeff Hammerbacher).]
  US-Non-Profit (\url{https://openathena.ai}), das mit akademischen
  Laboren Foundation Models baut, Engineers anstellt, Compute
  vermittelt und Multi-Universitäts-Kooperationen koordiniert. Größtes
  Projekt Marin (\url{https://marin.community}): offen entwickeltes
  LLM, zehn Vollzeit-Engineers.  Das Pretraining eines 535B-MoE mit
  23B aktiven Parametern auf 18T Tokens ist seit dem 5.\,11.\,2026
  abgeschlossen und war öffentlich verfolgbar.  Gespräche mit
  akademischen Laboren in den USA und Europa über Mathematik-Modelle
  laufen.  Frühere Läufe rechneten auf TPUs aus Googles TPU Research
  Cloud, die 768 GB200 für den großen Lauf stellte die Huang Foundation
  bereit.  Als Partner für Rechenzeit und Infrastruktur nennt Open
  Athena außerdem Lambda, CoreWeave und AWS, als Geldgeber Stiftungen
  wie Schmidt Sciences und die Sloan Foundation.  Eine vergleichbare
  Organisationsschicht und Unterstützung dieser Art gibt es in
  Deutschland bislang nicht.  Im Oktober 2026 hat Open Athena
  beschrieben, warum für das Post-Training ein offener Datenfundus
  fehlt (Abschnitt~\ref{sec:dsp}).
\item[Project Numina (Pietro Pellegrino).]
  Autoformalisierung mit Lean-4-Agenten (\enquote{Numina Fuse}),
  öffentliche Formalisierung der dreidimensionalen Kakeya-Vermutung,
  NuminaMath-Datensatz.
\item[Prover- und Formalisierungsmodelle.]
  Goedel-Prover-V2-8B und DeepSeek-Prover-7B (Beweise),
  StepFun-Formalizer-7B (Autoformalisierung); Plattform
  \url{https://prove2.me} für kollaborative Lean-Formalisierung.
\end{description}

\subsection{Öffentliche Modellprojekte in Europa}

Größen, Rechenzeit und Budgets stehen in Tabelle~\ref{tab:kosten}.

\begin{description}
\item[OpenGPT-X und Teuken-7B.] Fraunhofer IAIS und neun Partner,
  BMWK-gefördert 2022--2025, trainiert auf JUWELS Booster
  (Abschnitt~\ref{sec:dsp}).
\item[EuroLLM.] 9B-Modell auf 400 H100 (MareNostrum~5), EuroHPC- und
  Horizon-Europe-finanziert.
\item[Apertus (Swiss AI Initiative).] ETH, EPFL und CSCS; 8B und 70B,
  trainiert auf Alps. Aktuell ist Apertus~1.5.  Apertus~2.0 ist
  angekündigt, aber noch nicht erschienen.
\item[OpenEuroLLM.] Seit Februar 2025, 20 Partner unter Leitung der
  Karls-Universität Prag und von AMD Silo AI, 37~Mio.\ EUR (laut
  Presse), über 10~Mio.\ GPU-Stunden auf LUMI, Leonardo, JUPITER und
  MareNostrum~5. Das für Sommer 2026 angekündigte erste Modell gibt es
  bislang nur als Zwischen-Checkpoints (\enquote{Prelude 9B}, rund 10T
  Tokens, Apache~2.0), ohne Instruction-Tuning und ohne
  veröffentlichten Vergleich mit anderen Modellen; ein 32B-Modell ist
  im Training.
\item[hessian.AI, LeoLM, Occiglot.] Continued Pretraining im Bereich
  Tausender GPU-Stunden auf dem Cluster \enquote{42}.
\end{description}

\subsection{Öffentliche Recheninfrastruktur}

\begin{description}
\item[JUPITER und GCS.] Erstes europäisches Exascale-System am JSC
  Jülich, 500~Mio.\ EUR, rund \num{24000} GH200. Die Hälfte der
  Rechenzeit vergibt das Gauss Centre for Supercomputing an deutsche
  Projekte (Abschnitt~\ref{sec:kosten}).
\item[AI Factories.] Die 19 EuroHPC AI Factories sind das
  gegenwärtige KI-Angebot für die Wissenschaft; die beiden deutschen
  Standorte sind HammerHAI (Stuttgart, 55~Mio.\ EUR, etwa 860
  Blackwell-GPUs, zweites Halbjahr 2026) und JAIF (Jülich,
  55~Mio.\ EUR).
\item[AI Gigafactories.] Die EU schreibt bis zu sieben
  Rechenzentren der Industrieklasse mit jeweils \num{100000} oder mehr
  KI-Beschleunigern aus, finanziert mit bis zu 10~Mrd.\ EUR öffentlich
  (hälftig EU und Standort-Mitgliedstaat), die mindestens
  20~Mrd.\ EUR private Investitionen auslösen sollen. Die
  ursprüngliche Angebotsfrist (12.\,11.\,2026) gilt nicht mehr.  Seit
  dem 5.\,11.\,2026 endet sie für Lot~2 am 3.\,12.\,2026 und für
  Lot~1 am 4.\,2.\,2027, weil Mitgliedstaaten ihre Haushaltsverfahren
  noch abschließen müssen.  Der Betrieb soll 18~Monate nach der
  Auswahl beginnen, die bisher für Anfang 2027 vorgesehen war. EU und
  18 Mitgliedstaaten kaufen von den Betreibern Rechenzeit ein
  (gemeinsame Beschaffung, darunter Deutschland); die Ausschreibung
  verlangt von den Konsortien einen langfristig garantierten Anteil
  dieser Rechenzeit für die EU und den Standortmitgliedstaat. Diese
  Zeit kann laut Verordnung (EU) 2026/150 für kostenfreien Zugang
  verwendet werden, also auch für Forschung; welcher Anteil davon der
  Wissenschaft tatsächlich zufällt, ist offen und politisch zu
  verhandeln. Die Gigafactories sind die Industrie-Ebene: Die
  Wissenschaft hat mit JUPITER, GCS und den AI Factories einen
  eigenen, vorgelagerten Zugang. Berichteter deutscher Bieterverbund
  ist Telekom/Schwarz in Lübbenau (bis zu \num{100000} GPUs, ca.\
  11~Mrd.\ EUR).  Das ist Presseberichterstattung, kein bestätigter
  Stand.  SAP hat sich zurückgezogen.\footnote{Quellen in
  Abschnitt~\ref{sec:quellen}, Stand 17.\,11.\,2026.}
\item[GWDG und KISSKI.] Die GWDG, Rechenzentrum der Universität
  Göttingen und der Max-Planck-Gesellschaft, betreibt mit dem
  KI-Servicezentrum KISSKI (seit November 2022 vom BMFTR gefördert)
  den Dienst Chat AI mit einer OpenAI-kompatiblen API. Zugang über die
  Academic Cloud für Angehörige von Hochschulen und
  Forschungseinrichtungen. Die offenen Modelle, darunter GLM-5.3-Flash,
  DeepSeek~V4 Flash, Qwen3.5-397B und Qwen3.8-27B, laufen im eigenen
  Rechenzentrum, sind kostenlos und speichern laut GWDG keine
  Eingaben. Modelle von OpenAI und Anthropic werden über Microsoft
  Azure vermittelt und sind nicht kostenlos. Dazu kommen RAG,
  ein Coding-Assistent, Hosting und Feintuning eigener Modelle sowie
  GPU-Rechenzeit. Damit gibt es für Horizont~1 bereits ein Angebot für
  die deutsche Wissenschaft; zur Größe der dort betriebenen Modelle
  siehe Abschnitt~\ref{sec:querschnitt}.
\end{description}

\subsection{Hosting und Werkzeuge}

\begin{description}
\item[EU-Inferenzanbieter.] Cortecs (\url{https://cortecs.ai/}),
  Scaleway, OVHcloud, IONOS, STACKIT, Nebius (Finnland); Preise in
  Abschnitt~\ref{sec:kosten}.
\item[Agenten-Harnesses.] Claude Code, Codex (proprietär); OpenCode
  (\url{https://opencode.ai}) als offene Alternative
  (Abschnitt~\ref{sec:agenten}).
\end{description}

%=====================================================================
\section{Glossar}
\label{sec:glossar}

\begin{description}[style=sameline, itemsep=0.15em]
\item[Agent, agentisch] Ein LLM, das in einer Schleife Werkzeuge
  aufruft (Dateien lesen, Code ausführen, suchen) und Zwischenergebnisse
  weiterverarbeitet, statt nur eine Antwort zu erzeugen.
\item[Aktive Parameter] Bei MoE-Modellen der Teil der Parameter, der
  pro Token tatsächlich gerechnet wird. Bestimmt Geschwindigkeit und
  Trainingsaufwand; die Gesamtzahl bestimmt den Speicherbedarf. Bei
  dichten Modellen sind alle Parameter aktiv
  (Abschnitt~\ref{sec:groessen}).
\item[Autoformalisierung] Übersetzung informeller mathematischer
  Aussagen und Beweise in eine formale Sprache wie Lean, so dass ein
  Kernel sie prüfen kann. Die verbleibende Schwachstelle ist die
  Entsprechung von formaler und informeller Aussage.
\item[B, T] Milliarden (\emph{billion}) bzw.\ Billionen
  (\emph{trillion}) Parameter oder Tokens.
\item[Destillation] Ein kleineres Modell lernt aus den Ausgaben eines
  größeren. Um Größenordnungen billiger als Reinforcement Learning.
\item[FLOP, $6NT$] Gleitkommaoperationen. Der Trainingsaufwand ist
  näherungsweise sechsmal (aktive) Parameter mal Tokens.
\item[Frontiermodell] Die jeweils leistungsfähigsten Modelle der
  großen US-Firmen, derzeit Claude Fable~5.2 (Anthropic), GPT-6.2
  Astra (OpenAI) und Gemini~4 Carbon (Google), Stand 17.\,11.\,2026.
\item[GPU-Stunde] Eine Stunde Rechenzeit auf einer GPU; Standardmaß
  für Trainingskosten. H100: etwa 3~EUR Miete, 1,3~kWh Energie.
\item[Harness] Die Software um das Modell herum: Prompts, Werkzeuge,
  Kontextverwaltung, Schleifensteuerung. Entscheidet wesentlich über
  die agentische Leistung.
\item[HBM, NVLink] Schneller Speicher direkt auf der GPU bzw.\ die
  schnelle Verbindung zwischen GPUs eines Knotens.
\item[Hyperscaler] Betreiber sehr großer Rechenzentren, die Rechenzeit
  oder Modellbetrieb als Dienstleistung anbieten.
\item[Kernel (Lean)] Der kleine, vertrauenswürdige Prüfer, der formale
  Beweise verifiziert. Ein vom Kernel akzeptierter Beweis ist korrekt;
  ob die formale Aussage der gemeinten entspricht, prüft er nicht.
\item[KV-Cache] Zwischenspeicher für den Kontext jedes Nutzers; wächst
  linear mit der Kontextlänge und kann mehr Speicher brauchen als die
  Gewichte.
\item[MFU] Model FLOPs Utilization: Anteil der theoretischen
  Rechenleistung, der im Training tatsächlich genutzt wird, typisch
  30--45~Prozent.
\item[MoE] Mixture of Experts: Architektur, bei der pro Token nur
  einige von vielen Teilnetzen (Experten) aktiv sind.
\item[Open Weights, Open Source] Open Weights: die Gewichte sind
  herunterladbar. Open Source im engeren Sinn: auch Daten, Rezept und
  Code sind offen. Die meisten \enquote{offenen} Modelle sind nur
  Open Weights.
\item[Post-Training, SFT, RL] Die Phasen nach dem Pretraining:
  überwachtes Feintuning auf Beispielen (SFT) und Reinforcement
  Learning mit Belohnungen, etwa durch verifizierbare Ergebnisse.
\item[Pretraining] Das Grundtraining auf Billionen Tokens, der
  teuerste Rechenschritt.
\item[PUE] Power Usage Effectiveness: Gesamtstrom des Rechenzentrums
  geteilt durch IT-Strom. 1,2 ist gut, deutscher Durchschnitt 1,55.
\item[Quantisierung] Speichern der Gewichte mit weniger Bits (4 oder
  8 statt 16). Halbiert bzw.\ viertelt den Speicherbedarf bei geringem
  Qualitätsverlust.
\item[RAG] Retrieval-Augmented Generation: Das Modell erhält passende
  Dokumente aus einer Datenbank in den Kontext, statt sie auswendig zu
  wissen.
\item[Souveränität] Hier: Betrieb oder Training ohne Abhängigkeit von
  US-Anbietern, mit Daten, die in der eigenen Organisation oder
  zumindest in der EU bleiben.
\item[TDM, \S\,60d UrhG] Text- und Data-Mining; die Schranke, die
  Forschungseinrichtungen das Training auf urheberrechtlich
  geschützten Texten erlaubt.
\item[Token] Die Einheit, in der Modelle Text verarbeiten, grob drei
  Viertel eines Wortes.
\end{description}

%=====================================================================
\section{Quellen}
\label{sec:quellen}

Abgerufen im Oktober und November 2026. Rechenwege und der Prüfstatus
einzelner Angaben (berichtet oder verifiziert) sind in den
Recherchenotizen zu diesem Bericht festgehalten.

{\small
\subsection*{Modelle (Abschnitt~\ref{sec:groessen})}
\begin{itemize}[itemsep=0.1em]
  \item Qwen3.8-27B: \url{https://huggingface.co/Qwen/Qwen3.8-27B}
  \item Nemotron~3 Super 120B-A12B:
    \url{https://build.nvidia.com/nvidia/nemotron-3-super-120b-a12b/modelcard}
  \item GLM-5.3-Flash: \url{https://huggingface.co/zai-org/GLM-5.3-Flash}
  \item GLM-5.3: \url{https://huggingface.co/zai-org/GLM-5.3}
  \item DeepSeek-V4.1-Flash:
    \url{https://huggingface.co/deepseek-ai/DeepSeek-V4.1-Flash}
  \item Kimi~K3: \url{https://huggingface.co/moonshotai/Kimi-K3}
  \item Artificial Analysis, Intelligence Index v4.3.2 und agentische
    Benchmarks (Terminal-Bench~4.0, GDPval-AA, AutomationBench-AA),
    Stand 16.\,10.\,2026: GLM-5.3-Flash
    \url{https://artificialanalysis.ai/models/glm-5-3-flash}, Claude
    Opus~4.8 \url{https://artificialanalysis.ai/models/claude-opus-4-8},
    Rangliste \url{https://artificialanalysis.ai/leaderboards/models}
  \item Artificial Analysis, Intelligence Index v5.0.1, Stand
    17.\,11.\,2026: Rangliste wie oben, Methodik
    \url{https://artificialanalysis.ai/methodology/intelligence-benchmarking},
    Kimi~K3.1 Pro
    \url{https://artificialanalysis.ai/models/kimi-k3-1-pro}
\end{itemize}

\subsection*{Selbst hosten oder einkaufen (Abschnitt~\ref{sec:hosten}, Stand 17.\,11.\,2026)}
\begin{itemize}[itemsep=0.1em]
  \item Cortecs, Modellkatalog und Preisbildung:
    \url{https://cortecs.ai/serverlessModels},
    \url{https://cortecs.ai/pricing}
  \item OpenRouter, Endpunkte für GLM-5.3-Flash:
    \url{https://openrouter.ai/z-ai/glm-5.3-flash}
  \item Server mit acht H200, Preis und Leistungsaufnahme (UpStation,
    Stand 03.\,11.\,2026): \url{https://ai.upstation.io/gpu-servers}
  \item vLLM-Rezept für GLM-5.3-Flash (Speicherbedarf, Hardware):
    \url{https://recipes.vllm.ai/zai-org/GLM-5.3-Flash}
  \item Durchsatzmessung auf acht H200 (Projekt Kiln):
    \url{https://github.com/foxl-ai/kiln/blob/f5a7faf3bf/docs/price-performance.md}
\end{itemize}

\subsection*{Trainingskosten (Abschnitt~\ref{sec:kosten}, Tabelle~\ref{tab:kosten})}
\begin{itemize}[itemsep=0.1em]
  \item DeepSeek-V3 Technical Report, Dezember 2024:
    \url{https://arxiv.org/html/2412.19437v1}
  \item DeepSeek-R1, RL-Kosten (The Register, 19.\,09.\,2025):
    \url{https://www.theregister.com/2025/09/19/deepseek_cost_train/}
  \item Llama~3.1 405B, Modellkarte:
    \url{https://huggingface.co/meta-llama/Llama-3.1-405B}
  \item Llama~4 Maverick, Modellkarte:
    \url{https://huggingface.co/meta-llama/Llama-4-Maverick-17B-128E-Instruct}
  \item Llama~3, Paper (MFU, Unterbrechungen):
    \url{https://arxiv.org/abs/2407.21783}
  \item gpt-oss, Modellkarte: \url{https://arxiv.org/html/2508.10925v1}
  \item Olmo~3, Paper: \url{https://arxiv.org/html/2512.13961}
  \item BLOOM, Paper: \url{https://arxiv.org/pdf/2211.05100}
  \item Nemotron-4 340B, Report: \url{https://arxiv.org/pdf/2406.11704}
  \item Nemotron~3 Super, Report: \url{https://arxiv.org/pdf/2604.12374}
  \item Nemotron~3 Super, Pretraining-Dokumentation:
    \url{https://docs.nvidia.com/nemotron/0.1.0/nemotron/super3/pretrain.html}
  \item Qwen3 Technical Report: \url{https://arxiv.org/html/2505.09388v1}
  \item Qwen2.5 Technical Report: \url{https://arxiv.org/html/2412.15115v1}
  \item Apertus, Report: \url{https://arxiv.org/pdf/2509.14233}
  \item Teuken-7B, Fraunhofer IAIS:
    \url{https://www.iais.fraunhofer.de/en/press-events/press-releases/press-release-241126.html}
  \item EuroLLM, Report: \url{https://arxiv.org/html/2506.04079v1}
  \item Epoch AI, Grok~4 Training Resources:
    \url{https://epoch.ai/data-insights/grok-4-training-resources}
  \item Epoch AI, How much does it cost to train frontier AI models:
    \url{https://epoch.ai/blog/how-much-does-it-cost-to-train-frontier-ai-models}
  \item Cottier et al., Kostenanteile Personal und Energie:
    \url{https://arxiv.org/abs/2405.21015}
  \item Stanford AI Index 2026:
    \url{https://hai.stanford.edu/news/inside-the-ai-index-12-takeaways-from-the-2026-report}
\end{itemize}

\subsection*{Daten und Recht}
\begin{itemize}[itemsep=0.1em]
  \item FineWeb, Paper: \url{https://arxiv.org/html/2406.17557}
  \item FineWeb-2:
    \url{https://huggingface.co/datasets/HuggingFaceFW/fineweb-2}
  \item DCLM:
    \url{https://huggingface.co/datasets/mlfoundations/dclm-baseline-1.0}
  \item Nemotron-Pretraining-Dataset:
    \url{https://huggingface.co/datasets/nvidia/Nemotron-Pretraining-Dataset-sample}
  \item German Commons: \url{https://arxiv.org/abs/2510.13996}
  \item \S\,60d und \S\,44b UrhG:
    \url{https://www.gesetze-im-internet.de/urhg/__60d.html},
    \url{https://www.gesetze-im-internet.de/urhg/__44b.html}
  \item OLG Hamburg, LAION (Hogan Lovells, Dezember 2025):
    \url{https://www.insidetechlaw.com/blog/2025/12/machine-readable-opt-outs-and-ai-training-hamburg-court-clarifies-copyright-exceptions}
  \item LG München~I, GEMA gegen OpenAI, 42~O~14139/24 (LTO; Bird \& Bird):
    \url{https://www.lto.de/recht/nachrichten/n/42o1413924-lg-muenchen-i-gema-openai-chatgpt-songtexte-ki},
    \url{https://cm.twobirds.com/en/insights/2025/landmark-ruling-of-the-munich-regional-court-(gema-v-openai)-on-copyright-and-ai-training}
  \item AI Act, Art.~2, 3 und 53, Anhang~XI:
    \url{https://artificialintelligenceact.eu/article/2/},
    \url{https://artificialintelligenceact.eu/article/3/},
    \url{https://artificialintelligenceact.eu/article/53/},
    \url{https://artificialintelligenceact.eu/annex/11/}
  \item Europäische Kommission, Leitlinien zu den Pflichten der
    Anbieter von Modellen mit allgemeinem Verwendungszweck
    (18.\,07.\,2025), Rn.~17 und~51:
    \url{https://digital-strategy.ec.europa.eu/en/library/guidelines-scope-obligations-providers-general-purpose-ai-models-under-ai-act}
  \item Open Athena, \emph{There Is No Post-Training Data Commons}
    (15.\,10.\,2026):
    \url{https://openathena.ai/blog/there-is-no-posttraining-data-commons/}
  \item DSGVO, Art.~28 (Auftragsverarbeitung):
    \url{https://eur-lex.europa.eu/legal-content/DE/TXT/?uri=CELEX:32016R0679}
  \item CLOUD Act, 18 U.S.C. \S\,2713:
    \url{https://www.law.cornell.edu/uscode/text/18/2713}
\end{itemize}

\subsection*{Hardware, Rechenzeit und Strom}
\begin{itemize}[itemsep=0.1em]
  \item JUPITER (FZJ): \url{https://www.fz-juelich.de/en/jupiter}
  \item GCS, JUPITER-Factsheet:
    \url{https://jards.gauss-centre.eu/PublicFiles/FactSheets/FactSheet-GCS-NIC-JUPITER.pdf}
  \item EuroHPC, Large-Scale Access AI Factories:
    \url{https://www.eurohpc-ju.europa.eu/large-scale-access-ai-factories_en}
  \item HammerHAI (HLRS):
    \url{https://www.hlrs.de/news/detail/hammerhai-a-new-supercomputer-for-artificial-intelligence}
  \item OpenEuroLLM, Rechenzeit (AI Sweden):
    \url{https://www.ai.se/en/news/openeurollm-secures-strategic-access-multiple-supercomputers-milestone-european-language-models}
  \item H100-Kaufpreise (CloudZero):
    \url{https://www.cloudzero.com/blog/h100-gpu-cost/}
  \item Europäische GPU-Clouds (Qovery):
    \url{https://www.qovery.com/blog/european-gpu-cloud-alternatives-aws-azure-comparison}
  \item Energie von OLMo~2 (Morrison et al., ICLR 2025):
    \url{https://arxiv.org/pdf/2503.05804}
  \item BDEW-Strompreisanalyse 08/2026:
    \url{https://www.bdew.de/media/documents/BDEW_Strompreisanalyse_082026.pdf}
  \item Eurostat, Strompreise:
    \url{https://ec.europa.eu/eurostat/statistics-explained/index.php?title=Electricity_price_statistics}
  \item \S\,3, \S\,11 und \S\,12 EnEfG:
    \url{https://www.gesetze-im-internet.de/enefg/__3.html},
    \url{https://www.gesetze-im-internet.de/enefg/__11.html},
    \url{https://www.gesetze-im-internet.de/enefg/__12.html}
\end{itemize}

\subsection*{AI Gigafactories (Stand 13.\,10.\,2026, Fristen 17.\,11.\,2026)}
\begin{itemize}[itemsep=0.1em]
  \item EuroHPC JU, Start der Ausschreibung:
    \url{https://www.eurohpc-ju.europa.eu/eurohpc-joint-undertaking-launches-ai-gigafactories-call-2026-07-30_en}
  \item EuroHPC JU, Ausschreibung (mit Fristverlängerung vom
    05.\,11.\,2026):
    \url{https://www.eurohpc-ju.europa.eu/call-tenders-selection-artificial-intelligence-gigafactory-consortia-and-establishment-ai_en}
  \item Europäische Kommission, IP/26/1708:
    \url{https://ec.europa.eu/commission/presscorner/detail/en/ip_26_1708}
  \item Verordnung (EU) 2026/150:
    \url{https://eur-lex.europa.eu/legal-content/EN/TXT/?uri=CELEX%3A32026R0150}
  \item Deutscher Bieterverbund, unbestätigt (heise, 01.\,12.\,2025):
    \url{https://heise.de/-11098104}
  \item SAP nicht beteiligt (dpa via Marketscreener):
    \url{https://in.marketscreener.com/quote/stock/SAP-SE-436555/news/No-joint-German-bid-for-AI-gigafactory-50284021/}
\end{itemize}

\subsection*{Initiativen, Hosting und Werkzeuge (Abschnitt~\ref{sec:projekte})}
\begin{itemize}[itemsep=0.1em]
  \item SAIR Open Math Model Initiative (Tao, 18.\,09.\,2026):
    \url{https://terrytao.wordpress.com/2026/09/18/sairs-open-math-model-initiative/},
    \url{https://sair.foundation/open-math-model/}
  \item Open Athena, Partner und Geldgeber:
    \url{https://openathena.ai/about/}
  \item Marin, Unterstützer und Abschluss des Pretrainings (Stand
    17.\,11.\,2026):
    \url{https://marin.community},
    \url{https://github.com/marin-community/marin},
    \url{https://marin.community/blog/2026/11/09/marin-535b-architecture/}
  \item prove2.me: \url{https://prove2.me}
  \item OpenEuroLLM, Prelude 9B (Stand 17.\,11.\,2026):
    \url{https://openeurollm.eu/prelude},
    \url{https://huggingface.co/openeurollm/prelude}; Projektstart und
    Partner: \url{https://openeurollm.eu/blog/launch-press-release};
    Budget (TechCrunch, 16.\,02.\,2025):
    \url{https://techcrunch.com/2025/02/16/open-source-llms-hit-europes-digital-sovereignty-roadmap}
  \item GWDG/KISSKI, Chat AI und Modellliste (Stand 16.\,10.\,2026):\newline
    \url{https://docs.hpc.gwdg.de/services/ai-services/chat-ai/index.html},
    \url{https://docs.hpc.gwdg.de/services/ai-services/chat-ai/models/index.html},
    \url{https://kisski.gwdg.de/leistungen/services/}
  \item Cortecs: \url{https://cortecs.ai/}
  \item OpenCode: \url{https://opencode.ai}
  \item Epoch AI: \url{https://epoch.ai}
\end{itemize}
}

\end{document}
